\chapter{Introdução}
\label{sec:introducao}

Este capítulo contextualiza o problema investigado e apresenta a motivação para
o desenvolvimento do LEAFEON. A discussão parte das características do ensino
de Arquitetura de Software, examina as dificuldades associadas à realização e
ao acompanhamento de atividades práticas e situa a proposta em relação ao
EEVEE, plataforma educacional adotada como ponto de partida para este trabalho.
Ao final, são apresentados a questão de pesquisa, os objetivos, o percurso
metodológico e a organização do texto.

\section{Motivação e Contexto}
\label{sec:motivacao-contexto}

A arquitetura de software descreve estruturas de um sistema, seus elementos e
as relações estabelecidas entre eles. Também expressa decisões que condicionam
atributos de qualidade, entre os quais desempenho, segurança, disponibilidade e
modificabilidade. Seu estudo, portanto, não se limita à memorização de padrões e
definições. O estudante precisa interpretar um contexto, reconhecer restrições,
comparar alternativas e justificar decisões cujos efeitos se tornam visíveis ao
longo do desenvolvimento e da evolução do sistema
\cite{bass_software_2012,pantoja_training_2024}.

Essas características tornam o ensino da disciplina particularmente difícil.
Parte dos conceitos arquiteturais é abstrata, e seus efeitos nem sempre podem
ser observados em exercícios pequenos ou isolados. Por outro lado, a adoção de
projetos mais realistas introduz dificuldades de escopo, infraestrutura,
trabalho em equipe e acompanhamento. Uma revisão de 56 experiências
educacionais identificou a natureza abstrata e ambígua da arquitetura, a falta
de experiência prática proporcionada por projetos simples e a necessidade de
ir além de métodos tradicionais entre os problemas mais relatados na literatura
\cite{galster_teaching_2016,pantoja_training_2024}.

No curso de Bacharelado em Sistemas de Informação considerado nesta pesquisa,
o programa da disciplina reúne tópicos de naturezas distintas. São estudados
estilos e organizações arquiteturais, como microsserviços, monólitos modulares,
arquitetura hexagonal, sistemas orientados a eventos e ambientes multitenant;
princípios e métodos, como SOLID, DDD, Event Storming, Clean Code e padrões de
projeto; e tecnologias de integração, como GraphQL, gRPC, API Gateway e
mensageria, além de segurança em APIs. O desafio pedagógico não está apenas na
quantidade de assuntos, mas em permitir que o estudante relacione tecnologia,
decisão de projeto e atributo de qualidade. Essa combinação entre fundamentos,
padrões, princípios, análise e avaliação também aparece nos cursos examinados
pela literatura \cite{pantoja_training_2024}.

A aprendizagem ativa oferece um caminho para tratar esse problema, pois coloca
o estudante na posição de analisar situações, formular soluções e discutir as
consequências de suas escolhas. Uma meta-análise de 225 estudos realizados em
cursos de ciências, tecnologia, engenharia e matemática encontrou aumento de
desempenho e redução das taxas de reprovação nas turmas que adotaram estratégias
ativas em comparação com aulas predominantemente expositivas
\cite{freeman_active_2014}. Especificamente no ensino de arquitetura de
software, uma atividade na qual os estudantes representavam componentes e
simulavam suas interações favoreceu a compreensão de estilos arquiteturais
\cite{georgas_supporting_2013}. Esses resultados não eliminam a importância da
fundamentação teórica, mas indicam que ela pode ser articulada a situações nas
quais o estudante aplica e discute os conceitos apresentados.

A adoção de atividades frequentes, contudo, aumenta a quantidade de produções
que precisam ser acompanhadas. Esse aspecto é relevante porque o retorno perde
parte de sua utilidade quando chega somente depois de encerrado o processo de
resolução. Ferramentas de avaliação automatizada permitem aplicar critérios
repetíveis e devolver informações a cada tentativa. As revisões da área,
entretanto, mostram que essas ferramentas se concentram na correção funcional
de programas, em geral por meio de testes, e oferecem cobertura limitada de
aspectos como legibilidade, manutenibilidade e qualidade de projeto
\cite{keuning_feedback_2018,messer_grading_2024}. Não se pressupõe, portanto,
que toda decisão arquitetural possa ser corrigida automaticamente. A automação
é tratada neste trabalho como apoio ao estudante e ao professor, e não como
substituta da avaliação docente.

O caráter prático da proposta também se relaciona à aproximação entre formação
acadêmica e atuação profissional. Estudos sobre educação em Engenharia de
Software relatam diferenças entre as competências desenvolvidas na universidade
e aquelas esperadas pela indústria, sobretudo quanto à aplicação do
conhecimento em situações reais, à comunicação, ao trabalho em equipe e ao
pensamento crítico \cite{garousi_closing_2020}. Projetos de desenvolvimento,
projetos de código aberto, estudos de caso e aprendizagem baseada em problemas
são algumas das estratégias empregadas para aproximar o ensino de arquitetura
das condições encontradas no desenvolvimento de sistemas reais
\cite{pantoja_training_2024}.

Embora diferentes estratégias tenham sido experimentadas, a literatura não
aponta uma solução única para os problemas do ensino de arquitetura. Há uma
lacuna na combinação de estratégias pedagógicas, e ainda são poucas as
experiências relatadas com ambientes tutores on-line voltados à disciplina
\cite{pantoja_training_2024}. Esse resultado não demonstra, por si só, a
inexistência de ferramentas; indica, de maneira mais restrita, que permanece
aberta a investigação de ambientes capazes de organizar atividades práticas,
executar experimentos de maneira controlada e registrar evidências do processo
de realização dessas atividades.

O EEVEE oferece uma base concreta para essa investigação. Nele, o professor
pode preparar atividades de programação e testes automatizados, e o estudante
pode desenvolver uma solução em um ambiente Web, submetê-la e consultar o
resultado da execução. A plataforma registra tentativas e reaplica os mesmos
critérios de teste a diferentes submissões, características que também permitem
seu uso em experimentos educacionais \cite{eevee_repository_2026}. Seu modelo
atual, porém, foi construído principalmente em torno de atividades de
programação. Para utilizá-lo no ensino de arquitetura, será necessário
investigar como representar cenários, executar configurações arquiteturais,
observar seus efeitos e fornecer retorno sobre decisões que podem admitir mais
de uma solução aceitável.

Diante desse contexto, este trabalho propõe o LEAFEON, uma extensão da
plataforma EEVEE voltada à experimentação e à aprendizagem de Arquitetura de
Software. A proposta aproveita a base tecnológica e o fluxo educacional do
EEVEE e os amplia com recursos para representar cenários, executar atividades
práticas e observar consequências de decisões arquiteturais. Os requisitos e a
forma específica de integração serão definidos a partir das necessidades
pedagógicas e técnicas identificadas durante a pesquisa.

Com base no problema exposto, estabelece-se a seguinte questão de pesquisa:
\textbf{como uma plataforma interativa de experimentação pode apoiar a
realização de atividades práticas de Arquitetura de Software no ensino
superior?}

\section{Objetivos}
\label{sec:objetivos}

\textbf{Objetivo geral.} Desenvolver e avaliar o LEAFEON como uma extensão do
EEVEE destinada a apoiar a realização de atividades práticas e a compreensão de
conceitos de Arquitetura de Software no ensino superior.

\textbf{Objetivos específicos.} Para alcançar o objetivo geral, foram definidos
os seguintes objetivos específicos:

\begin{enumerate}
    \item caracterizar, a partir da literatura e do contexto da disciplina, os
    desafios associados ao ensino prático de Arquitetura de Software e as
    estratégias utilizadas para enfrentá-los;

    \item identificar os requisitos pedagógicos e técnicos necessários para a
    realização, o acompanhamento e a avaliação de atividades práticas de
    Arquitetura de Software;

    \item projetar e desenvolver o LEAFEON, definindo sua forma de integração
    com o EEVEE e implementando os recursos necessários para a execução de
    atividades de experimentação arquitetural;

    \item elaborar atividades práticas que permitam aos estudantes aplicar,
    observar e comparar conceitos e decisões de Arquitetura de Software;

    \item aplicar o LEAFEON em uma turma da disciplina de Arquitetura de
    Software, registrando dados de uso, resultados das atividades e percepções
    dos participantes; e

    \item analisar os resultados da aplicação para identificar contribuições,
    limitações e oportunidades de melhoria no LEAFEON.
\end{enumerate}

\section{Metodologia}
\label{sec:metodologia}

\subsection{Delineamento da pesquisa}

Este trabalho será conduzido segundo a Design Science Research (DSR). Essa
abordagem é adequada a pesquisas que procuram compreender um problema e, ao
mesmo tempo, projetar e avaliar um artefato destinado a intervir no contexto
estudado. Na DSR, o desenvolvimento do artefato não constitui um fim isolado:
sua construção deve ser orientada pelo problema, fundamentada no conhecimento
existente e acompanhada de avaliação sistemática
\cite{hevner_design_2004,dresch_design_2015}.

O artefato investigado será o LEAFEON, compreendendo a plataforma, sua
integração com o EEVEE, as atividades educacionais implementadas e os mecanismos
de acompanhamento utilizados na avaliação. O processo de pesquisa foi organizado
com base nas etapas propostas para a DSR: identificação do problema, definição
dos objetivos da solução, projeto e desenvolvimento, demonstração, avaliação e
comunicação \cite{peffers_methodology_2007}. Essas etapas não serão tratadas
como uma sequência estritamente linear. Os resultados das avaliações poderão
levar à revisão dos requisitos, das atividades e da implementação.

\subsection{Etapas e entregáveis}

O percurso metodológico será composto pelas seguintes etapas:

\begin{enumerate}
    \item \textbf{Busca na literatura:} serão pesquisados estudos sobre ensino
    de Arquitetura de Software, aprendizagem ativa, ambientes educacionais e
    avaliação automatizada. A busca será documentada por meio das bases
    consultadas, strings de busca e critérios de seleção. O entregável desta
    etapa será o conjunto inicial de estudos candidatos e o protocolo utilizado
    para recuperá-los;

    \item \textbf{Seleção e análise dos estudos:} os trabalhos recuperados serão
    avaliados inicialmente pelo título e resumo e, quando considerados
    pertinentes, por leitura integral. Serão extraídos dados sobre problemas
    abordados, estratégias pedagógicas, ferramentas, formas de avaliação e
    resultados. Os entregáveis serão o portfólio de estudos selecionados e a
    síntese dos trabalhos correlatos;

    \item \textbf{Definição dos requisitos e da avaliação:} os achados da
    literatura serão confrontados com o programa da disciplina e com as
    necessidades do professor. Como resultado, serão produzidos a especificação
    dos requisitos pedagógicos e técnicos, a definição inicial das atividades
    práticas e o plano de avaliação do artefato;

    \item \textbf{Projeto e desenvolvimento:} o LEAFEON será desenvolvido de
    forma iterativa. Cada ciclo compreenderá a seleção de requisitos, a
    implementação de um incremento e sua verificação técnica e pedagógica. Os
    entregáveis incluirão os modelos da solução, o código-fonte, versões
    executáveis e atividades de experimentação arquitetural;

    \item \textbf{Demonstração em sala de aula:} uma versão considerada estável
    será utilizada em uma turma de Arquitetura de Software. O entregável desta
    etapa será o conjunto de dados produzidos durante a aplicação, incluindo
    resultados das atividades, tentativas, tempos de execução e respostas aos
    instrumentos de pesquisa;

    \item \textbf{Avaliação do artefato:} os dados coletados serão analisados
    para verificar o atendimento aos requisitos e investigar as contribuições
    do LEAFEON para a realização de atividades práticas e para a aprendizagem.
    O entregável será um relatório de avaliação com os resultados quantitativos
    e qualitativos, as limitações observadas e as melhorias recomendadas; e

    \item \textbf{Refinamento e comunicação:} os resultados da avaliação serão
    utilizados para priorizar correções e melhorias. Quando houver tempo e
    viabilidade, uma nova versão retornará à etapa de demonstração e avaliação.
    Os entregáveis finais serão a versão avaliada do LEAFEON, este trabalho e a
    comunicação dos resultados em publicação acadêmica.
\end{enumerate}

O fluxo metodológico será representado com a Business Process Model and
Notation (BPMN), de modo a explicitar a sequência das atividades, os entregáveis
e o retorno da avaliação para um novo ciclo de desenvolvimento. A notação será
utilizada conforme a especificação BPMN 2.0.2 \cite{omg_bpmn_2014}.

\subsection{Planejamento da avaliação}

A avaliação será conduzida com estudantes de uma turma de Arquitetura de
Software após a realização de atividades práticas no LEAFEON. Será adotado um
estudo de uso e aceitação com um único grupo de participantes, sem divisão entre
grupo experimental e grupo de controle. Essa escolha mantém a avaliação
compatível com o caráter exploratório do trabalho e com a aplicação da
ferramenta em uma turma, sem atribuir à pesquisa a finalidade de demonstrar uma
relação causal entre o uso do LEAFEON e o desempenho acadêmico.

Durante a aplicação, serão coletados registros produzidos pela plataforma,
como resultados das atividades, tempo de conclusão, número de tentativas e taxa
de conclusão. Após o uso, será aplicado um questionário baseado no Modelo de
Aceitação de Tecnologia, contemplando facilidade de uso percebida, utilidade
percebida e intenção de uso \cite{davis_perceived_1989}. As respostas objetivas
serão registradas em escala Likert de cinco pontos, e questões abertas permitirão
que os participantes relatem dificuldades, contribuições percebidas e sugestões
de melhoria.

Os dados quantitativos serão analisados de forma descritiva, por meio de
frequências, medidas de tendência central e sínteses dos registros de uso. As
respostas abertas serão organizadas por temas recorrentes. A interpretação
conjunta dessas evidências buscará verificar a viabilidade de uso do LEAFEON, a
aceitação da ferramenta e sua contribuição percebida para a realização das
atividades práticas. Dessa forma, os resultados permitirão avaliar o artefato e
orientar seu refinamento, sem serem apresentados como comprovação definitiva de
melhoria da aprendizagem.

A participação na pesquisa deverá ser voluntária, com esclarecimento sobre os
dados coletados e sua finalidade. Os dados utilizados na análise serão
anonimizados, e a participação ou recusa não deverá produzir prejuízo acadêmico.
Antes da aplicação, serão observados os procedimentos éticos exigidos pela
instituição para pesquisas realizadas com estudantes.
